% Propietario conservado: /Users/ruben/Documents/New project/output/REVISION_ARGUMENTAL_APERTURAS_CIERRES_20260915/III/spanish_source/manuscrito/sections/06_pell_barbero.tex
\section{Evaluación en la escala de Pell y funcional constitutivo de Barbero--Immirzi}
\label{iii:iii:sec:pell-barbero}

La relación entre la respuesta del vacío y el momento de Catalán permite
componer dos niveles del transporte: la profundidad ponderada del carácter
de cuarto de giro y la evaluación de los canales angulares. La escala de
Pell distingue un argumento de la familia de momentos; los canales del
vacío reciben, en cambio, los argumentos determinados por el par angular.
El funcional de Barbero--Immirzi se expresa sobre estos últimos mediante
la misma respuesta constitutiva y su inversa. El vínculo se establece así
entre operaciones definidas, conservando la posición de cada evaluación
en la construcción APP--TRIT--TPK.

\subsection{Momento orientado y evaluación en la unidad de Pell}

Escribamos
\begin{equation}
 G_{\rm Cat}:=\mathcal C_2(1)
 =\sum_{k\geq0}\frac{(-1)^k}{(2k+1)^2},\qquad
 \lambda_{\rm P}=2+\sqrt3,\qquad
 \rho_{\rm P}=2-\sqrt3=\lambda_{\rm P}^{-1}.
 \label{iii:iii:eq:pell-scales}
\end{equation}
La constante $G_{\rm Cat}$ es el valor extremo del momento construido en
\eqref{iii:iii:eq:c2}. El par $\lambda_{\rm P},\rho_{\rm P}$ corresponde a
las orientaciones expansiva y contractiva de la unidad de Pell:
$(2+\sqrt3)(2-\sqrt3)=1$. La identidad de norma $2^2-3\cdot1^2=1$
fija la reciprocidad. La especialización del momento en $\rho_{\rm P}$
es posterior a la construcción de su carácter orientado.

La ecuación diferencial \eqref{iii:iii:eq:c2-integral}, con condición nula
en el origen, da
\begin{equation}
 \mathcal D\mathcal C_2(s)=\arctan s,\qquad
 \mathcal C_2(s)=\int_0^s\frac{\arctan t}{t}\,dt,
 \qquad \mathcal D=s\frac{d}{ds}.
 \label{iii:iii:eq:pell-moment-integral}
\end{equation}
En efecto, integrar $\mathcal D^2\mathcal C_2=s/(1+s^2)$ respecto de
$ds/s$ produce $\arctan s$; una segunda integración produce la expresión
indicada. El integrando se prolonga continuamente en $t=0$ con valor $1$.
La forma angular de esta integral permite determinar exactamente la
evaluación en $\rho_{\rm P}$.

\begin{lemma}[Integral angular del carácter impar]
\label{iii:iii:lem:log-tangent}
Para $0<\theta\leq\pi/4$,
\begin{equation}
 -\int_0^\theta\log(\tan x)\,dx
 =\sum_{\substack{n\geq1\\n\text{ impar}}}
      \frac{\sin(2n\theta)}{n^2}.
 \label{iii:iii:eq:log-tangent}
\end{equation}
\end{lemma}
\begin{proof}
Para $0<a<1$, la expansión del logaritmo dentro del disco unidad implica
\begin{align*}
 L_a(x)
 &:=\log|1-ae^{2ix}|-\log|1+ae^{2ix}|\\
 &=-2\sum_{\substack{n\geq1\\n\text{ impar}}}
       \frac{a^n\cos(2nx)}{n}.
\end{align*}
La serie converge uniformemente en $x$ para cada $a<1$; su integración
término a término proporciona
\[
 -\int_0^\theta L_a(x)\,dx
 =\sum_{\substack{n\geq1\\n\text{ impar}}}
       \frac{a^n\sin(2n\theta)}{n^2}.
\]
Cuando $a\uparrow1$, el cociente de los módulos converge a
$|1-e^{2ix}|/|1+e^{2ix}|=\tan x$ para $0<x\leq\theta$.
La singularidad del logaritmo en el origen es integrable. Para justificar
el límite, puede restringirse $a\geq1/2$ y observar que
\[
 |1-ae^{2ix}|^2=(1-a)^2+4a\sin^2x
 \geq2\sin^2x,
 \qquad 1\leq|1+ae^{2ix}|\leq2
\]
en $0\leq x\leq\pi/4$. Junto a la cota superior
$|1-ae^{2ix}|\leq2$, estas desigualdades dan
$|L_a(x)|\leq c+|\log x|$ con una constante independiente de $a$.
La convergencia dominada permite pasar al límite en la integral.
En la serie integrada se aplica la dominación por $\sum n^{-2}$.
Los dos límites prueban \eqref{iii:iii:eq:log-tangent}.
\end{proof}

\begin{theorem}[Evaluación Catalán--Pell]
\label{iii:iii:thm:catalan-pell}
El momento orientado satisface
\begin{align}
 \mathcal C_2(\rho_{\rm P})
 &=\frac23G_{\rm Cat}-\frac{\pi}{12}\log\lambda_{\rm P},
 \label{iii:iii:eq:catalan-pell}\\
 \mathcal D\mathcal C_2(\rho_{\rm P})&=\frac{\pi}{12},&
 \mathcal D^2\mathcal C_2(\rho_{\rm P})&=\frac14.
 \label{iii:iii:eq:catalan-pell-derivatives}
\end{align}
\end{theorem}
\begin{proof}
La sustitución $t=\tan x$ en \eqref{iii:iii:eq:pell-moment-integral},
seguida de integración por partes, da
\begin{equation}
 \mathcal C_2(\tan\theta)
 =\theta\log(\tan\theta)
   -\int_0^\theta\log(\tan x)\,dx.
 \label{iii:iii:eq:angular-moment}
\end{equation}
El término de frontera en cero se anula porque $x\log(\tan x)\to0$.
La fórmula de la tangente de una diferencia proporciona
\[
 \tan\frac{\pi}{12}
 =\tan\left(\frac\pi4-\frac\pi6\right)
 =\frac{\sqrt3-1}{\sqrt3+1}=2-\sqrt3.
\]
Queda por evaluar la serie de \eqref{iii:iii:eq:log-tangent} en
$\theta=\pi/12$. Sea $\chi_{-4}(n)=(-1)^{(n-1)/2}$ para $n$ impar.
En las seis clases impares módulo $12$ se verifica
\begin{equation}
 2\sin\frac{n\pi}{6}
 =\chi_{-4}(n)+3\,\mathbf1_{3\mid n}\chi_{-4}(n/3).
 \label{iii:iii:eq:character-pell}
\end{equation}
La comprobación finita es
\[
\begin{array}{c|rrrrrr}
 n\bmod12&1&3&5&7&9&11\\ \hline
 2\sin(n\pi/6)&1&2&1&-1&-2&-1\\
 \chi_{-4}(n)&1&-1&1&-1&1&-1\\
 3\mathbf1_{3\mid n}\chi_{-4}(n/3)&0&3&0&0&-3&0
\end{array}
\]
y ambos miembros son periódicos módulo $12$ sobre los índices impares.
La convergencia absoluta permite separar los sumandos e introducir $n=3m$
en el segundo término:
\begin{align*}
 \sum_{\substack{n\geq1\\n\text{ impar}}}
       \frac{\sin(n\pi/6)}{n^2}
 &=\frac12\sum_{n\text{ impar}}\frac{\chi_{-4}(n)}{n^2}
  +\frac32\sum_{m\text{ impar}}\frac{\chi_{-4}(m)}{(3m)^2}\\
 &=\left(\frac12+\frac16\right)G_{\rm Cat}
 =\frac23G_{\rm Cat}.
\end{align*}
Las sumas sin límite inferior en esta expresión recorren los enteros
positivos. Sustituyendo en \eqref{iii:iii:eq:angular-moment} y usando
$\log\rho_{\rm P}=-\log\lambda_{\rm P}$ se obtiene
\eqref{iii:iii:eq:catalan-pell}. La primera derivada resulta de
$\arctan\rho_{\rm P}=\pi/12$; la segunda, de
\[
 \frac{\rho_{\rm P}}{1+\rho_{\rm P}^2}
 =\frac1{\rho_{\rm P}+\rho_{\rm P}^{-1}}=\frac14.
\]
\end{proof}

El coeficiente $2/3$ procede de la descomposición del carácter en
\eqref{iii:iii:eq:character-pell}: la restricción a los múltiplos de $3$
modifica el peso de profundidad por el factor $3^{-2}$. La evaluación
reúne, por tanto, el cuarto de giro y la subdivisión dodecafásica en una
identidad exacta de momentos. Los coeficientes se fijan antes de cualquier
comparación decimal.

\subsection{Composición de fase y profundidad ponderada}

La representación de profundidad precisa qué conserva esta evaluación.
En $\ell^2(\mathbb N_{>0})$, definamos
\[
 N\delta_n=n\delta_n,\qquad U_4\delta_n=i^n\delta_n,
 \qquad Q_4=\frac{U_4-U_4^*}{2i}.
\]
Para $0<s<1$, el operador $N^{-2}s^NQ_4$ es de clase traza, y
\begin{equation}
 \mathcal C_2(s)=\operatorname{Tr}(N^{-2}s^NQ_4).
 \label{iii:iii:eq:pell-depth-trace}
\end{equation}
Esta igualdad se obtiene evaluando sobre la base $\delta_n$; las
entradas diagonales son $s^n\sin(n\pi/2)/n^2$. El intercambio
$U_4\leftrightarrow U_4^*$ invierte $Q_4$ y, con él, el momento
orientado. El índice $n$ y la ponderación $s^n$ permanecen fijos.

La composición se realiza antes de tomar la traza. Para
$a\in\mathbb Z/4\mathbb Z$, sea $B(s,a)=s^NU_4^a$. Como los dos
factores son diagonales en la misma base,
\begin{equation}
 B(s,a)B(t,b)=B(st,a+b),\qquad
 B(s,1)^4=s^{4N}.
 \label{iii:iii:eq:pell-depth-composition}
\end{equation}
El retorno de fase conserva la contracción acumulada. En particular,
$B(\rho_{\rm P},1)^4=\rho_{\rm P}^{4N}$, cuyo autovalor en profundidad
$n$ es $\rho_{\rm P}^{4n}$. Esta realización da un significado
operatorio a la evaluación contractiva; la suma escalar del momento
agrega las profundidades después de conservar su índice y orientación.
Los argumentos constitutivos $s_\pm=q_\pm^{30}$ siguen determinados por
el par angular de \eqref{iii:iii:eq:recover-angular}. Las fórmulas anteriores
no contienen una igualdad entre esos argumentos y $\rho_{\rm P}$.

\subsection{Incidencia dodecafásica y funcional de retorno}

La contribución de Barbero--Immirzi utiliza el transporte angular y la
incidencia hexada--octada. Sean $\mathsf H\subset\mathsf O$ una hexada
y una octada con el par marcado de la hexada conservado. Los conjuntos de
banderas pertinentes son
\begin{equation}
 \mathcal F_6=\mathsf H\times\binom{\mathsf H}{2}
 \lhook\joinrel\longrightarrow
 \mathcal F_8=\mathsf O\times\binom{\mathsf H}{2},
 \qquad |\mathcal F_6|=90,\quad |\mathcal F_8|=120.
 \label{iii:iii:eq:barbero-flags}
\end{equation}
Los cardinales son $6\binom62$ y $8\binom62$; la misma incidencia
conserva el defecto normalizado
$(90-120)/(24\binom62)=-1/12$. Esta cuenta identifica los dos sectores
empleados por el retorno, mientras el calendario determina sus
multiplicidades relativas.

Para $m\in\{90,120\}$, el retorno trítico sobre un canal contractivo es
\begin{equation}
 S_m(q)=\sum_{j\geq0}q^{m+3mj}
       =\frac{q^m}{1-q^{3m}},\qquad 0<q<1.
 \label{iii:iii:eq:barbero-return-series}
\end{equation}
Cada término corresponde a la altura inicial $m$ y a $j$ retornos de
longitud $3m$. La serie geométrica conserva, por tanto, tanto la altura
inicial como el período del retorno. El calendario dodecafásico
$\mathbb Z/12\mathbb Z$ aporta doce sectores y una única incidencia
orientada de frontera, denotada $\partial_{\rm ret}$. En el módulo
libre de sucesos, su cadena fundamental es
\begin{equation}
 c_{12}^{+}
 =\sum_{j\in\mathbb Z/12\mathbb Z}[j,\mathcal F_6]
  -[\partial_{\rm ret},\mathcal F_8].
 \label{iii:iii:eq:barbero-cycle}
\end{equation}
El lector lineal $\mathscr L_q$ asigna $S_{90}(q)$ a cada generador
$[j,\mathcal F_6]$ y $S_{120}(q)$ al generador de frontera. Se obtiene
\begin{equation}
 \mathscr L_q(c_{12}^{+})=12S_{90}(q)-S_{120}(q).
 \label{iii:iii:eq:barbero-weights}
\end{equation}
Los pesos $12$ y $-1$ son las multiplicidades orientadas de la cadena.
La orientación conjugada satisface $c_{12}^{-}=-c_{12}^{+}$, con los
mismos tipos de incidencia.

\begin{proposition}[Coeficiente singular del retorno fundamental]
\label{iii:iii:prop:barbero-pole}
El coeficiente de $(1-q)^{-1}$ de
$\mathscr L_q(c_{12}^{+})$ es $1/24$.
\end{proposition}
\begin{proof}
Para cada entero $m>0$, la factorización
$1-q^{3m}=(1-q)\sum_{j=0}^{3m-1}q^j$ da
\[
 \lim_{q\uparrow1}(1-q)S_m(q)=\frac1{3m}.
\]
Aplicar esta identidad al funcional ya determinado en
\eqref{iii:iii:eq:barbero-weights} produce
\begin{equation}
 \lim_{q\uparrow1}(1-q)\mathscr L_q(c_{12}^{+})
 =\frac{12}{270}-\frac1{360}=\frac1{24}.
 \label{iii:iii:eq:barbero-pole}
\end{equation}
En la coordenada compleja $q-1$, el residuo es $-1/24$.
\end{proof}

La evaluación singular verifica una consecuencia de la incidencia y del
retorno trítico. El orden demostrativo queda fijado por
\eqref{iii:iii:eq:barbero-flags}--\eqref{iii:iii:eq:barbero-weights}; el
coeficiente singular se calcula después de construir la cadena.

\subsection{Factorización a través del operador constitutivo}

Conservemos la carta angular levantada de la
sección~\ref{iii:iii:sec:hojas}. Para hacer explícita la conversión angular,
escribamos
\begin{equation}
 A_{\deg}=\frac{180x}{\pi},\qquad
 C^*_{\deg}=\frac{180y}{\pi},\qquad
 q_\pm=e^{-x\mp y},\qquad s_\pm=q_\pm^{30}.
 \label{iii:iii:eq:barbero-angular-lift}
\end{equation}
La orientación considerada en \eqref{iii:iii:eq:rchannels} corresponde a
$0<y<x$; el intercambio de hojas se expresa por $y\mapsto-y$ en el
dominio más amplio $x>|y|$. Las exponenciales reales actúan sobre estos
levantamientos, que conservan los valores de $x\pm y$ durante la
composición.

El funcional orientado de Barbero--Immirzi es
\begin{equation}
 \gamma=\frac{\pi A_{\deg}C^*_{\deg}}{180}
   +12\bigl[S_{90}(q_-)-S_{90}(q_+)\bigr]
   -\bigl[S_{120}(q_-)-S_{120}(q_+)\bigr].
 \label{iii:iii:eq:barbero-functional}
\end{equation}
La primera componente es la evaluación angular; la segunda es la
diferencia del lector de la cadena fundamental en los canales conjugados.
Los argumentos se reciben del transporte angular previamente construido.
La fórmula determina un escalar adimensional sobre esa estructura de
partida.

Sea $\psi:(3/4,1)\to(0,1)$ la inversa de $f$ construida mediante
\eqref{iii:iii:eq:cubic}. Definamos
\begin{equation}
 \mathcal H(s)=\frac{12s^3}{1-s^9}
             -\frac{s^4}{1-s^{12}}
             -\frac{(\log s)^2}{20\pi},\qquad 0<s<1.
 \label{iii:iii:eq:barbero-h}
\end{equation}
Esta función es real analítica en su dominio. Los exponentes $3,4,9,12$
resultan de expresar las alturas $90,120$ y los retornos $270,360$ en
unidades de $30$; el coeficiente de la parte logarítmica se fija al
transportar la contribución angular.

\begin{theorem}[Expresión constitutiva orientada de Barbero--Immirzi]
\label{iii:iii:thm:barbero-factorization}
En la carta \eqref{iii:iii:eq:barbero-angular-lift},
\begin{equation}
 \gamma=\mathcal H(s_-)-\mathcal H(s_+)
       =\mathcal H\bigl(\psi(r_-)\bigr)
        -\mathcal H\bigl(\psi(r_+)\bigr).
 \label{iii:iii:eq:barbero-factorization}
\end{equation}
Sobre la representación bidimensional de canales,
\begin{equation}
 \gamma=-\operatorname{Tr}
    \left[\mathcal R\,\mathcal H\bigl(\psi(\mathcal V)\bigr)\right],
 \label{iii:iii:eq:barbero-trace}
\end{equation}
donde $\operatorname{Tr}$ es la traza ordinaria, con
$\operatorname{Tr}I=2$.
\end{theorem}
\begin{proof}
La sustitución $s=q^{30}$ en \eqref{iii:iii:eq:barbero-return-series} da
\[
 S_{90}(q)=\frac{s^3}{1-s^9},\qquad
 S_{120}(q)=\frac{s^4}{1-s^{12}}.
\]
Además, \eqref{iii:iii:eq:barbero-angular-lift} implica
\[
 \log s_\pm=-30(x\pm y)
 =-\frac\pi6(A_{\deg}\pm C^*_{\deg}).
\]
Por diferencia de cuadrados,
\begin{align}
 \frac{(\log s_+)^2-(\log s_-)^2}{20\pi}
 &=\frac{900\bigl[(x+y)^2-(x-y)^2\bigr]}{20\pi}\\
 &=\frac{180xy}{\pi}
 =\frac{\pi A_{\deg}C^*_{\deg}}{180}.
 \label{iii:iii:eq:barbero-log-term}
\end{align}
La suma de estas identidades prueba la primera igualdad de
\eqref{iii:iii:eq:barbero-factorization}. La segunda utiliza la inversión
$\psi(f(s_\pm))=s_\pm$ demostrada en la
sección~\ref{iii:iii:sec:constitutiva}.

Para la expresión operatoria, los proyectores de hoja tienen rango uno y
$\mathcal RP_\pm=\pm P_\pm$. El cálculo espectral proporciona
\[
 \mathcal H\bigl(\psi(\mathcal V)\bigr)
 =\mathcal H(s_+)P_++\mathcal H(s_-)P_-.
\]
Multiplicar por $\mathcal R$ y tomar la traza ordinaria da
$\mathcal H(s_+)-\mathcal H(s_-)$. Su opuesto es $\gamma$.
\end{proof}

La marca de orientación $\mathcal R$ forma parte del funcional. Si se
mantiene fija esa marca y se intercambian los canales de $\mathcal V$,
cambia el signo de $\gamma$. La evaluación del espectro sin sus etiquetas conserva
el par no ordenado, mientras \eqref{iii:iii:eq:barbero-trace} conserva la
contribución orientada. Si se emplea la traza normalizada
$\tau_2=\operatorname{Tr}/2$, la misma fórmula se escribe
$\gamma=-2\tau_2[\mathcal R\mathcal H(\psi(\mathcal V))]$.

En cambio, un cambio simultáneo de representación
$(\mathcal R,\mathcal V)\mapsto
(U\mathcal RU^{-1},U\mathcal VU^{-1})$ conserva el funcional:
el cálculo espectral se transporta por conjugación y la traza es
invariante. Intercambiar la respuesta respecto de una orientación fija
y transportar conjuntamente respuesta y orientación son, por tanto,
operaciones diferentes.

Esta normalización es compatible con las amplificaciones de
\eqref{iii:iii:eq:logvac}. Para
$\mathcal V_m=\mathcal V\otimes I_{9^m}$ y
$\mathcal R_m=\mathcal R\otimes I_{9^m}$,
\begin{equation}
 \gamma=-\frac1{9^m}\operatorname{Tr}
 \left[\mathcal R_m\mathcal H\bigl(\psi(\mathcal V_m)\bigr)\right].
 \label{iii:iii:eq:barbero-amplification}
\end{equation}
En efecto, el cálculo funcional conserva el producto tensorial y la traza
de $I_{9^m}$ es $9^m$. La multiplicidad de la representación se separa
de la evaluación orientada.

\subsection{Reconstrucción a partir de las coordenadas del vacío}

Las identidades \eqref{iii:iii:eq:four-identities} permiten expresar el mismo
funcional mediante la impedancia y la velocidad normalizadas:
\begin{align}
 \gamma={}&
 \mathcal H\!\left(\psi\!\left(
     \sqrt{\frac{\widehat Z}{\widehat c}}\right)\right)
 -\mathcal H\!\left(\psi\!\left(
     \frac1{\sqrt{\widehat Z\widehat c}}\right)\right).
 \label{iii:iii:eq:barbero-zc}
\end{align}
El dominio positivo de esta reconstrucción se describe por
\begin{equation}
 1<\widehat Z\widehat c<\frac{16}{9},\qquad
 \frac9{16}<\frac{\widehat Z}{\widehat c}<1.
 \label{iii:iii:eq:barbero-zc-domain}
\end{equation}
En efecto, esas desigualdades equivalen a
$3/4<r_\pm<1$ al sustituir las raíces positivas de
\eqref{iii:iii:eq:four-identities}. Para la orientación $y>0$ se conserva
además $r_-<r_+$, equivalentemente $\widehat Z<1$. Los puntos del
dominio se utilizan con la compatibilidad angular ya establecida por
el TPK; la inversión es una reconstrucción posterior de sus canales.
Las coordenadas son las internas de \eqref{iii:iii:eq:four}: un cambio de
unidades debe deshacerse mediante \eqref{iii:iii:eq:undo-units} antes de
aplicar $\psi$.

El intercambio $y\mapsto-y$ actúa como
$(\widehat Z,\widehat c)\mapsto(\widehat Z^{-1},\widehat c)$ y
produce $\gamma\mapsto-\gamma$. En el caso balanceado $y=0$,
ambos canales coinciden y \eqref{iii:iii:eq:barbero-factorization} da
$\gamma=0$. La velocidad normalizada conserva el producto de los
canales; su lectura conjunta con la impedancia recupera las dos
respuestas y la orientación necesaria para evaluar el funcional.
En ese caso balanceado $\mathcal V=rI$ tiene espectro degenerado:
la marca $\mathcal R$ se conserva como dato de representación, aunque
el escalar $r$ ya no distinga sus proyectores. Para espectro simple,
en cambio,
$P_+=(\mathcal V-r_-I)/(r_+-r_-)$ y $P_-=I-P_+$.

El resultado central de esta composición es la factorización
\eqref{iii:iii:eq:barbero-factorization}. El momento de Catalán aporta el
coeficiente orientado que determina $f$; la inversión cúbica recupera
el argumento de cada canal; la incidencia dodecafásica y el par angular
determinan $\mathcal H$. La escala de Pell distingue una evaluación
exacta de la familia, con la prueba de
\eqref{iii:iii:eq:catalan-pell}. Cada relación conserva su argumento y su
operación: la constante extrema, el momento contractivo y el funcional
constitutivo constituyen evaluaciones enlazadas de estructuras
especificadas.
